\input{../tex-inputs/latex-leseansicht-vorspann}

\section[Arthur und Olga Schnitzler an Gustav Schwarzkopf, 22. 5. 1904]{L04371 Arthur und Olga Schnitzler an Gustav Schwarzkopf, 22. 5. 1904}
\nopagebreak\mylabel{L04371v}
\rehead{ }\normalsize\beginnumbering\briefempfaengerindex{Schwarzkopf, Gustav@\textsc{Schwarzkopf, Gustav}!zzzSchnitzler, Olga@\emph{von Olga Schnitzler}!1904-05-222@{22. 5. 1904}|(be}\briefempfaengerindex{Schwarzkopf, Gustav@\textsc{Schwarzkopf, Gustav}!zzzSchnitzler, Arthur@\emph{von Arthur Schnitzler}!1904-05-222@{22. 5. 1904}|(be}
\toendnotes[C]{\smallbreak\pagebreak[2]}
\correspDesc{Versand  durch Arthur Schnitzler, Olga Schnitzler am 22. 5. 1904 in Taormina
\newline{}Übermittlung  am 23. 5. 1904 in Taormina
\newline{}Zustellung  am 25. 5. 1904 in 1/1 Wien 1
\newline{}Erhalt  durch Gustav Schwarzkopf am 25. 5. 1904 in Tiefer Graben 23}\toendnotes[C]{\smallbreak}
\Standort{DLA, A:Schnitzler, HS.1985.1.1897.}
\physDesc{Bildpostkarte, 182 Zeichen
\newline{}Handschrift Arthur Schnitzler: Bleistift, deutsche Kurrent
\newline{}Handschrift Olga Schnitzler: Bleistift, deutsche Kurrent
\newline{}Versand: 1) Stempel: »\nobreak{}\oindex{Taormina@\textbf{Taormina}|pwk}Taormina
                                       (\textcolor{gray}{M}essina), 2\textcolor{gray}{3} \textcolor{gray}{5} 04\nobreak{}«.   2) Stempel: »\nobreak{}\oindex{I., Innere Stadt@\textbf{I., Innere Stadt}|pwk}1/1 Wien 1, 25 5 04, 8–9½V., Bestellt\nobreak{}«. }\pstart{}{\pb}\textsc{Austria}\oindex{Österreich@\textbf{Österreich}|pw}\pend{}\pstart{}\textsc{Hrn Gustav Schwarzkopf}\pend{}\pstart{}\textsc{Vienna}\oindex{Wien@\textbf{Wien}|pw}\pend{}\pstart{}I. Tiefer Graben 23\oindex{Wien@\textbf{Wien}!I., Innere Stadt@\textbf{I., Innere Stadt}!Tiefer Graben 23@\textbf{Tiefer Graben 23}|pw}.\pend{}{\bigskip}
\pstart
           {\pb}\textcolor{gray}{\textbf{Etna\oindex{Ätna@\textbf{Ätna}|pw} dal Giardino del Grand Hôtel S. Domenico\oindex{San Domenico Hotel@\textbf{San Domenico Hotel}|pw}}}\hfill \textcolor{gray}{\textbf{Taormina}}\oindex{Taormina@\textbf{Taormina}|pw}\pend
           \vspace{1em}
\pstart
           \raggedleft{}{\pb}22. Mai 1904\pend
           \vspace{0.5em}
\pstart
           Nur daſs es etwas wärmer iſt.\pend
           
\pstart
           Herzliche Grüße! Ihr{\\[\baselineskip]}\spacefill\mbox{A.}\pend
           \leftskip=0em{}\selectlanguage{ngerman}\vspace{1em}
\pstart
           \noindent{}{[}hs. O. Schnitzler:{]} Und viel duftendes Grün und{ }ſehr viele Blüthen!\pend
           
\pstart
           Liebe Grüße!{\\[\baselineskip]}\spacefill\mbox{O. S.}\pend
           \leftskip=0em{}\selectlanguage{ngerman}\endnumbering\briefempfaengerindex{Schwarzkopf, Gustav@\textsc{Schwarzkopf, Gustav}!zzzSchnitzler, Olga@\emph{von Olga Schnitzler}!1904-05-222@{22. 5. 1904}|)be}\briefempfaengerindex{Schwarzkopf, Gustav@\textsc{Schwarzkopf, Gustav}!zzzSchnitzler, Arthur@\emph{von Arthur Schnitzler}!1904-05-222@{22. 5. 1904}|)be}\mylabel{L04371h}
\begin{anhang}
\end{anhang}\newcommand{\dateiname}{L04371}\newcommand{\titel}{Arthur und Olga Schnitzler an Gustav Schwarzkopf, 22. 5. 1904}\newcommand{\editorInnen}{Herausgegeben von Jahnke, SelmaMüller, Martin Anton}\input{../tex-inputs/latex-leseansicht-abspann}
